% SHARD-ID: LB-04-CORNER-A
% SHARD-TITLE: Corner A: asymptotic symmetry as boundary-bond insertions
% SHARD-SOURCES: claims/CLAIMS.md (rows WI, A1, A2, G0, A2-orbit-r1,
%   G0-soft-r1, and the Corner-A annotation block with amendments (i)--(iv));
%   theory/corner-a.md; theory/corner-a-kinks.md; theory/corner-a-goldstone.md;
%   theory/corner-a-pitfalls.md; theory/verdicts/corner-a-r3.md;
%   theory/checks/corner_a_check.py (C0--C11); docs/framing.md (orientation only)
% SHARD-CLAIMS: WI, A1, A2, G0, A2-orbit-r1, G0-soft-r1

% --- shard-local semantic macros (rule 10 of the writing guide) -------------
\providecommand{\PGL}{\operatorname{PGL}}      % projective linear group
\providecommand{\PU}{\operatorname{PU}}        % projective unitary group
\providecommand{\GLie}{\operatorname{GL}}      % general linear group
\providecommand{\Ug}{\operatorname{U}}         % unitary group
\providecommand{\Hcoh}{\operatorname{H}}       % group cohomology
\providecommand{\bond}[2]{(#1\,|\,#2)}         % a lattice bond
\providecommand{\Gdiag}{G_{\mathrm{diag}}}
\providecommand{\Aeff}{\A_{\mathrm{eff}}}      % effective asymptotic symmetry group
\providecommand{\chargealg}{\mathfrak{a}_\alpha} % twisted group algebra
\providecommand{\bondcur}{\mathcal{J}}         % virtual bond current (D10(d))
\providecommand{\act}{\triangleright}          % action of a charge on a state

\section{Corner A: asymptotic symmetry as boundary-bond insertions}
\label{sec:corner-a}

\subsection{What Corner A is about}

A global on-site symmetry, applied to a \emph{finite} region of the chain, does
nothing at all in the interior of that region and everything at its two edges.
This is the whole of Corner A in one sentence, and on a matrix-product state it
is a theorem rather than a slogan: the symmetry is absorbed into the virtual
(bond) indices, the absorptions cancel pairwise along the region, and what
survives is a single virtual matrix sitting on each of the region's two boundary
bonds.  Everything else in this section --- the endpoint torsor, the projective
charge algebra, the superselected kink sectors, the lattice Noether identity ---
is a consequence of that cancellation, read in one register or another.

\paragraph*{Notation recalled.}
We work on the infinite spin chain with quasi-local algebra $\alg$ and its local
subalgebras $\alg_W$, $W$ finite, as fixed in \cref{sec:mps-setting}.  A vacuum
is described by an injective canonical-form tensor $A_\alpha=(A_\alpha^s)$,
$A_\alpha^s\in M_\chi(\C)$, whose transfer operator $E_\alpha$ has a unique
modulus-one eigenvalue, second-largest modulus $\lambda_E<1$, and injectivity
length $n_0$ (products of $n_0$ tensors span $M_\chi(\C)$).  Window vectors
$\ket{\psi_\Lambda(A;b_l,b_r)}$ carry boundary vectors $b_l,b_r\in\C^\chi$ and
may be \emph{decorated}: a matrix $M$ inserted on a bond $\bond{m}{m+1}$, or a
different tensor placed at some sites; $\omega_\alpha^{M@b}$ denotes the state of
the vacuum decorated by $M$ on the bond $b$, and $C_\partial$ the uniform norm
constant for one such insertion.  The on-site symmetry is a compact group $G$
acting by $u(g)$, with $U_R(g)=\prod_{x\in R}u_x(g)$ on a finite region $R$ and
$(\mathcal{U}(g)A)^s=\sum_{s'}u(g)_{ss'}A^{s'}$ on tensors.  The covariant vacuum
family $\{A_\alpha\}_{\alpha\in\Omega_{\mathrm{vac}}}$ obeys the intertwining
relation
\begin{equation}
  \mathcal{U}(g)A_\alpha \;=\; e^{i\theta_\alpha(g)}\,V_\alpha(g)^{-1}\,
  A_{g\cdot\alpha}\,V_\alpha(g),
  \label{eq:IT}
\end{equation}
with $V_\alpha(g)$ unitary and $\theta_\alpha(g)$ real, the pair unique up to
$V\mapsto e^{i\varphi}V$.  We write $H_\alpha$ for the stabiliser of $\alpha$,
$\mathfrak{h}_\alpha$ for its Lie algebra and $\mathfrak{m}_\alpha$ for a
complementary (broken) subspace.  Following the campaign's overload convention,
$\omega_\alpha$ with no argument is the vacuum \emph{state}, while
$\omega_\alpha(h,g)$ with two group arguments is the projective
\emph{two-cocycle} of \cref{eq:cocycle} below; the argument types never coincide.
The definitions themselves live in \cref{sec:mps-setting} and
\cref{sec:symmetry-noether} and are not restated here.

\paragraph*{Large versus small, on a lattice.}
The background viewpoint we keep (a symmetry counts as physical, or ``large'',
exactly when it acts on the state space with a well-defined charge, not because
of a fall-off slogan) has a sharp lattice translation.  A truncated symmetry
operation is genuinely non-trivial at infinity precisely when the virtual matrix
it deposits on the surviving bond is \emph{not} a multiple of the identity: by
\cref{thm:a1}(b) the half-infinite string is then implemented on states but by no
strongly convergent sequence of operators, and by \cref{thm:a1}(c) the state it
produces really is a new one.  If instead $V_\alpha(g)$ is scalar, the two
insertions cancel exactly, the string sequence is constant, and the operation is
invisible --- ``small''.  The discriminator is thus a finite-dimensional
linear-algebra condition on $V_\alpha(g)$, not a decay rate.  What plays the role
of the choice of function spaces at infinity is the class of profiles $f$
allowed in a modulated charge: summable-of-bounded-variation packets and
$\delta$-normalised plane waves behave differently, and \cref{thm:g0}(c) is
stated separately in each.

\paragraph*{The charge-algebra viewpoint.}
Read as a charge algebra, the content of Corner A is that the bond insertions
carry a \emph{projective} representation of $G$: on padded windows the operators
$\mathcal{V}_b(g)$ compose with the multiplier $e^{i\omega_\alpha(h,g)}$, whose
class $[\omega_\alpha]\in \Hcoh^2(H_\alpha,\Ug(1))$ is the lattice analogue of a
central extension of the charge algebra --- and is the one-dimensional
symmetry-protected-topological index of \cref{sec:spt}.  Two disciplines must be
kept.  First, the extension is group-cohomological, not Lie-algebraic: for
$\mathfrak{h}_\alpha$ compact semisimple Whitehead's second lemma forces the
infinitesimal cocycle to be a coboundary, so a nonzero torsion class
$[\omega_\alpha]$ can coexist with an entirely ordinary Lie bracket for the
edge generators (the AKLT edge is the standard example).  Second, the algebra is
so far shown to act on finite-window vectors and on states, \emph{not} on the
GNS Hilbert space of the vacuum; the missing step is the split property, recorded
as a load-bearing sketch in \cref{scope:corner-a-sketches}.

\subsection{The telescoping identity}

The following lemma is the algebraic engine.  It is what makes the phase in
\cref{eq:IT} removable, and it produces the cocycle whose class is the SPT index.

\begin{lemma}[Composition of intertwiners, the cocycle, and normal ordering]
\label{thm:lemma-it}
Assume the setting above.  For all $g,h\in G$:
\begin{enumerate}
\item[(i)] $\theta_\alpha(hg)=\theta_\alpha(g)+\theta_{g\cdot\alpha}(h)$ and
  \begin{equation}
    V_{g\cdot\alpha}(h)\,V_\alpha(g) \;=\; e^{i\omega_\alpha(h,g)}\,V_\alpha(hg);
    \label{eq:cocycle}
  \end{equation}
\item[(ii)] restricted to $H_\alpha$, $\theta_\alpha$ is a homomorphism
  $H_\alpha\to \Ug(1)$ and $\omega_\alpha$ is a $\Ug(1)$-valued two-cocycle whose
  class $[\omega_\alpha]\in \Hcoh^2(H_\alpha,\Ug(1))$ is independent of the phase
  convention for $V_\alpha$, of the matrix-product gauge $A^s\mapsto Y^{-1}A^sY$,
  and of blocking sites;
\item[(iii)] in the unbroken case $H_\alpha=G$, the \emph{normally ordered}
  operators $\check{u}_\alpha(g):=e^{-i\theta_\alpha(g)}u(g)$ form a unitary
  representation of $H_\alpha$ satisfying \cref{eq:IT} with phase $1$;
\item[(iv)] in the unbroken case, $V_\alpha(g)\,r\,V_\alpha(g)^\dagger=r$ for the
  fixed point $r$ of the transfer operator, and
  $E\bigl(V_\alpha(g)YV_\alpha(g)^\dagger\bigr)=V_\alpha(g)E(Y)V_\alpha(g)^\dagger$.
\end{enumerate}
\end{lemma}

\provenance{--- (supporting lemma of WI and A1; no separate claim row)}%
{\texttt{theory/corner-a.md}, step (1.2)}%
{\texttt{theory/checks/corner\_a\_check.py} C0}%
{\texttt{theory/verdicts/corner-a-r3.md}}

\smallskip\noindent\textit{Idea of proof.}
Apply \cref{eq:IT} twice and compare with its single application at $hg$: the
on-site action composes as a representation, so the two right-hand sides agree,
and the uniqueness of the pair $(\theta_\alpha,V_\alpha)$ up to a phase forces
(i).  Associativity of the triple product $V(g_1)V(g_2)V(g_3)$ is the cocycle
identity; rephasing $V_\alpha\mapsto e^{i\varphi}V_\alpha$ changes
$\omega_\alpha$ by a coboundary, a gauge change conjugates $V_\alpha$ and leaves
$\omega_\alpha$ alone, and blocking $n$ sites reproduces \cref{eq:IT} with the
same $V_\alpha$ and phase $n\theta_\alpha(g)$, which gives (ii).  Part (iv) is the
covariance of the transfer operator, from unitarity of $u(g)$ together with
uniqueness of the fixed point.\hfill$\square$

\begin{theorem}[Truncated symmetry telescopes to the boundary (WI)]
\label{thm:wi}
\statusProved\;
Let $R=[a,b]$ be a finite interval, $g\in G$, and $b_l,b_r\in\C^\chi$ arbitrary.
Assume \emph{either} (W1) a finite window $\Lambda\supseteq R$ together with edge-bond
insertions, \emph{or} (W2) a window $\Lambda\supseteq[a-1,b+1]$, in which case
both insertions land on interior bonds.  Then, exactly, as window vectors,
\begin{equation}
  U_R(g)\,\bigl|\psi_\Lambda(A_\alpha;b_l,b_r)\bigr\rangle
  \;=\; e^{\,i\abs{R}\theta_\alpha(g)}\,
        \bigl|\psi_\Lambda(T_R^{(g)};b_l,b_r)\bigr\rangle,
  \label{eq:WI}
\end{equation}
where the decoration $T_R^{(g)}$ carries $A_\alpha$ on $\Lambda\setminus R$,
$A_{g\cdot\alpha}$ on $R$, the bond insertion $V_\alpha(g)^{-1}$ on
$\partial_-R=\bond{a-1}{a}$ and $V_\alpha(g)$ on $\partial_+R=\bond{b}{b+1}$.
As \emph{states}, with \textbf{no phase}, for every $O\in\alg$,
\begin{equation}
  \omega_\alpha\bigl(U_R(g)^\dagger\,O\,U_R(g)\bigr)
  \;=\; \omega_\alpha\bigl[T_R^{(g)}\bigr](O).
  \label{eq:WI-state}
\end{equation}
In particular, in the unbroken case the interior of $R$ is unchanged and the
entire effect of a symmetry applied to a region is two virtual operators on the
region's two boundary bonds.
\end{theorem}

\begin{scope}
The window hypothesis is not cosmetic.  For $\Lambda=R$ with interior insertions
only, the two boundary matrices fall outside the expression and \cref{eq:WI} is
\emph{false}: the measured discrepancy for a $\chi=2$ Pauli tensor is $1.658$,
against $0.0$ once edge-bond insertions are admitted.  The extensive phase
$e^{i\abs{R}\theta_\alpha(g)}$ belongs to the \emph{vector} identity only; a
phase times a state is the same state, so \cref{eq:WI-state} carries none, and
asserting a limit of the \emph{operator} $U_{[x,y]}(g)$ when
$\theta_\alpha(g)\neq0$ is meaningless.  The orientation --- $V_\alpha(g)^{-1}$
on the left bond, $V_\alpha(g)$ on the right --- is fixed by the convention in
\cref{eq:IT}; it is the opposite of the assignment in the campaign brief, which
is a historical document and is not edited.  A $\Z_2$ test tensor with
$V=Z=V^{-1}$ cannot see the difference and must not be used to confirm it; a
$\Ug(1)$ test with $V(t)=e^{itZ/2}$ gives error $5.6\cdot10^{-17}$ for the
correct orientation and $0.267$ for the flipped one.  Finally, only
\cref{eq:IT} is used: injectivity is needed to \emph{obtain} \cref{eq:IT} and to
make $(\theta_\alpha,V_\alpha)$ unique, not for the algebra.
\end{scope}

\provenance{WI}%
{\texttt{theory/corner-a.md}, step (1.3)}%
{\texttt{theory/checks/corner\_a\_check.py} C1, C1b, C2, C2b}%
{\texttt{theory/verdicts/corner-a-r3.md}}

\smallskip\noindent\textit{Idea of proof.}
Expand $U_R(g)\ket{\psi_\Lambda}$ in the physical basis: the effect is to replace
$A_\alpha^{s_x}$ by $(\mathcal{U}(g)A_\alpha)^{s_x}$ at each $x\in R$, and by
\cref{eq:IT} each such replacement is
$e^{i\theta_\alpha(g)}V_\alpha(g)^{-1}A_{g\cdot\alpha}^{s_x}V_\alpha(g)$.  In the
ordered product over $x=a,\dots,b$ every interior pair
$V_\alpha(g)V_\alpha(g)^{-1}$ collapses to $\mathbf{1}$ --- this is the
telescoping --- leaving exactly one factor at each end and $\abs{R}=b-a+1$
copies of the phase.  Passing to states, the phase cancels between bra and ket
and the decoration differs from $A_\alpha$ on finitely many sites and
bonds.\hfill$\square$

\subsection{The unbroken case: endpoints, charge algebra, and the obstruction}

Throughout this subsection $H_\alpha=G$ and we work normally ordered, so
$\theta_\alpha\equiv0$.  Recall from \cref{sec:symmetry-noether} the objects
$N_\alpha:=\{g\in G: V_\alpha(g)\in\C^\times\mathbf{1}\}$ (a normal subgroup),
the induced map $\rho_\alpha:G\to\PGL(\chi)$, $\rho_\alpha(g)=[V_\alpha(g)]$, the
effective group $\Aeff:=G/N_\alpha\cong\rho_\alpha(G)$, the endpoint space
$E^\alpha_b=\{\omega_\alpha^{M@b}: M\in\GLie(\chi)\}$ at a bond $b$, the twisted
group algebra $\chargealg=\C_{\omega_\alpha}[G]$, and the space
$\mathcal{W}_{\Lambda,b}$ spanned by window vectors carrying one insertion on
$b$, with $\iota_{\Lambda,b}:M_\chi(\C)\to\mathcal{W}_{\Lambda,b}$ the map that
places the insertion.  A window is \emph{padded about $b$} if it contains at
least $n_0$ sites on each side of $b$ and $b_l,b_r\neq0$.

\begin{theorem}[Endpoint torsor, projective charge algebra, and the cohomological obstruction (A1)]
\label{thm:a1}
\statusProved\;
In the unbroken case, normally ordered, with every clause mentioning
$\mathcal{W}_{\Lambda,b}$ or $\mathcal{V}_b$ additionally assuming a window
padded about $b$:
\begin{enumerate}
\item[(a)] \emph{(charges live on one bond)} For finite $R$ the right insertion
  is invisible to observables to its left and the left insertion invisible to
  those to its right.  Hence the half-infinite operations exist \emph{as maps on
  states},
  \[
    1_{[x,\infty)}g \act \omega_\alpha
    \;=\; \omega_\alpha^{\,V_\alpha(g)^{-1}@\bond{x-1}{x}},
  \]
  and the convergence is \emph{exact} --- eventually constant on each $\alg_W$
  --- not merely asymptotic.
\item[(b)] \emph{(non-implementability, as an equivalence)} The sequence
  $\bigl(U_{[x,y]}(g)\Omega_A\bigr)_y$ fails to be Cauchy in $\Hilb_A$
  \textbf{if and only if} $V_\alpha(g)\notin\C^\times\mathbf{1}$.  For scalar
  $V_\alpha(g)$ the strings stabilise exactly.  Thus the half-infinite symmetry
  is implemented on states, but by no strongly convergent sequence of operators.
\item[(c)] \emph{(the endpoint bijection)}
  $\omega_\alpha^{M@b}=\omega_\alpha^{M'@b}$ \textbf{iff} $M'\in\C^\times M$;
  hence $[M]\mapsto\omega_\alpha^{M@b}$ is a canonical bijection and $E^\alpha_b$
  is a $\PGL(\chi)$-torsor.
\item[(d)] \emph{(the charge algebra, in the two registers it actually has)}
  \begin{enumerate}
  \item[(d1)] On a window padded about $b$, the map $\iota_{\Lambda,b}$ is
    injective, so
    $\mathcal{V}_b(M):=\iota_{\Lambda,b}\circ L_M\circ\iota_{\Lambda,b}^{-1}$ is a
    well-defined linear operator, and with $\mathcal{V}_b(g):=\mathcal{V}_b(V_\alpha(g))$
    \begin{equation}
      \mathcal{V}_b(h)\,\mathcal{V}_b(g)\;=\;e^{i\omega_\alpha(h,g)}\,\mathcal{V}_b(hg):
      \label{eq:twisted-rep}
    \end{equation}
    a \emph{linear} representation of $\chargealg=\C_{\omega_\alpha}[G]$ in which
    the multiplier acts nontrivially.  \emph{Padding is necessary, not
    cosmetic}: without it the kernel of $\iota_{\Lambda,b}$ need not be invariant
    under left multiplication and $\mathcal{V}_b$ is not an operator at all.
  \item[(d2)] The induced action on \emph{states} kills the multiplier and is the
    genuine homomorphism $\rho_\alpha:G\to\PGL(\chi)$ with
    $\ker\rho_\alpha=N_\alpha$; it is bond-independent and padding-independent.
  \item[(d3)] The window action (d1) \emph{is} a lift of $\rho_\alpha$, but a
    projective one, and such a lift exists for every class $[\omega_\alpha]$.
    What $[\omega_\alpha]$ obstructs is \textbf{removing the multiplier}: it is
    exactly the obstruction to lifting $\rho_\alpha:G\to\PGL(\chi)$ to an honest
    homomorphism $G\to\Ug(\chi)$, being the pullback under $\rho_\alpha$ of the
    extension class of $1\to\Ug(1)\to\Ug(\chi)\to\PU(\chi)\to1$.  It does
    \emph{not} obstruct the projective window action, which always exists.
  \end{enumerate}
\item[(e)] \emph{(stabiliser and orbit)} The stabiliser of $\omega_\alpha$ in
  $G_L\times G_R$ is
  \[
    S_\alpha=\{(g_L,g_R): g_Lg_R^{-1}\in N_\alpha\}\;\supseteq\;\Gdiag,
  \]
  with equality \textbf{iff} $N_\alpha=\{e\}$, and the orbit is in bijection with
  $\Aeff=G/N_\alpha$, a genuine group.  The coset space
  $\A=(G_L\times G_R)/\Gdiag$ is \emph{not} the orbit unless $N_\alpha=\{e\}$,
  and is a group only for abelian $G$.
\item[(f)] \emph{(what acts on what)} $\Aeff$ acts simply transitively on the
  orbit $\rho_\alpha(G)\subseteq E^\alpha_b$, and does not move the sector label
  $(\alpha,\alpha)$, which is trivial in the unbroken case.
\item[(g)] \emph{(invariance of the index)} $[\omega_\alpha]$ is constant along
  any continuous path of $G$-symmetric injective canonical-form tensors, provided
  $\Hcoh^2(G,\Ug(1))$ is discrete.
\end{enumerate}
\end{theorem}

\begin{scope}
The claims register carries this theorem with the scope qualifier ``padded-window
and state registers as stated'', and the clauses displayed in its headline row
are (a)--(e); (f) and (g) are proved in the same shard and are recorded here
because \cref{sec:spt} uses (g).  What is \emph{not} proved is the physical,
GNS-level realisation: that $\omega_\alpha^{M@b}$ is a normal state of
$\pi_\alpha(\alg)''$, equivalently that $\chargealg$ acts on the vacuum Hilbert
space rather than on finite-window spaces.  That is the split-property sketch of
\cref{scope:corner-a-sketches}, and it is load-bearing: the background
desideratum of a charge algebra with a well-defined action on the physical state
space is \emph{not} met by this theorem.  Reading (d) as an action on an edge
Hilbert space is therefore an over-read.  Completeness of $[\omega_\alpha]$ as a
classifier of one-dimensional SPT phases is cited from the literature, not proved
here.
\end{scope}

\provenance{A1}%
{\texttt{theory/corner-a.md}, step (1.4)}%
{\texttt{theory/checks/corner\_a\_check.py} C6, C8, C8b}%
{\texttt{theory/verdicts/corner-a-r3.md}}

\smallskip\noindent\textit{Idea of proof.}
Clause (a) is a transfer-operator computation on top of \cref{thm:wi}: in
canonical form the left environment is $\mathbf{1}$ and the right environment is
the fixed point $r$, and a unitary insertion $M$ replaces them by
$M^\dagger\mathbf{1}M=\mathbf{1}$ and $MrM^\dagger$; since
$V_\alpha(g)rV_\alpha(g)^\dagger=r$ by \cref{thm:lemma-it}(iv), the far insertion
simply disappears once the region's edge has passed the support of the
observable, which is why the limit is exact rather than asymptotic.

Clause (c) is the multilinear step.  Writing the reduced density matrix of a
window containing the bond, its matrix elements are
$\operatorname{tr}[P_1MP_2rQ_2^\dagger M^\dagger Q_1^\dagger]$ with the four words
$P_1,P_2,Q_1,Q_2$ varying \emph{independently}; by injectivity each ranges over a
spanning set of $M_\chi(\C)$, and both sides of the equality with the undecorated
state are multilinear, so the equality extends to all of $M_\chi(\C)^{\times4}$.
Specialising the words to matrix units turns this into $MZM^\dagger=\gamma Z$ for
every $Z$, whence $M\in\C^\times\mathbf{1}$.  Note the normalisation:  the
correct target is $\gamma Z$ with $\gamma>0$, not $Z$.

Clause (b) uses the overlaps $\braket{\Psi_y,\Psi_{y'}}=\rho_{y'-y}$ coming from
translation invariance.  Being Cauchy forces $\rho_n=1$ for \emph{every} $n\ge1$,
not merely $\rho_n\to1$; that in turn makes the one-sided string act trivially on
every local algebra, and (c) then forces $V_\alpha(g)$ scalar.  Conversely a
scalar $V_\alpha(g)=c\mathbf{1}$ gives boundary insertions $c^{-1}\mathbf{1}$ and
$c\mathbf{1}$ whose product is the identity, so the sequence is constant.

Clause (d1) rests on injectivity of $\iota_{\Lambda,b}$ on padded windows: the
coefficient $b_l^\dagger PMQb_r$ is linear in each of the two words, which span
$M_\chi(\C)$, and $b_l,b_r\neq0$, so $v^\dagger(M-M')w=0$ for all $v,w$.  Without
padding this fails concretely: for the injective tensor $A^0=\operatorname{diag}(1,2)$,
$A^1=X$ with $V=Z$, one site each side, and a suitable $N$, one has
$\iota(N)=0$ while $\norm{\iota(ZN)}_\infty=4$, so $\ker\iota$ is not invariant
and the rule ``$\iota(M)\mapsto\iota(NM)$'' defines nothing.  Given injectivity,
\cref{eq:twisted-rep} is \cref{eq:cocycle} transported through $\iota$.  Clause
(d2) is (c) again --- states see $M$ only modulo scalars --- and (d3) is the
observation that a lift with $\omega_\alpha\equiv0$ exists iff the class vanishes.

Clause (e) composes the two half-infinite residues: the left string leaves
$V_\alpha(g_L)$ on the bond and the right string leaves $V_\alpha(g_R)^{-1}$, so
the composite residue is $V_\alpha(g_L)V_\alpha(g_R)^{-1}$, trivial on states
exactly when $g_Lg_R^{-1}\in N_\alpha$; the map
$(g_L,g_R)\mapsto g_Lg_R^{-1}N_\alpha$ is then a bijection
$(G\times G)/S_\alpha\to G/N_\alpha$.  Clause (g) is continuity of an isolated
spectral projection along the path, giving a continuous map into a discrete
set.\hfill$\square$

\begin{remark}[Two instances that pin the orbit down]
A $\chi=1$ product vacuum with a $\Z_2$ symmetry has $V_\alpha\equiv1$, hence
$N_\alpha=G$, $S_\alpha=G\times G$, and a \emph{one-point} orbit.  The AKLT
vacuum with $G=\Z_2\times\Z_2$ has $V_\alpha\in\{\mathbf{1},X,Z,XZ\}$, hence
$N_\alpha=\{e\}$ and a four-point orbit of pairwise distinct states.  Any reading
of the orbit as ``a copy of the symmetry group'' fails on the first example.
\end{remark}

\subsection{The broken case: sector jumps and the classification of vacuum pairs}

Now let $\abs{\Omega_{\mathrm{vac}}}\ge2$, fix $\alpha$ and $g\in G$ with
$\beta:=g\cdot\alpha\neq\alpha$.  Recall the kink sectors $\Kk_{\alpha\beta}$ of
\cref{sec:symmetry-noether}: states with factorised vacuum asymptotics
$\omega_\alpha$ at $-\infty$ and $\omega_\beta$ at $+\infty$.  Sectors with
distinct labels are disjoint, and the label of a decorated matrix-product state
is read off from its tensors at $\pm\infty$ at exponential rate.

\begin{theorem}[A broken truncated symmetry creates a kink only in the limit (A2)]
\label{thm:a2}
\statusProved\;
For each fixed $g\notin H_\alpha$:
\begin{enumerate}
\item[(a)] \emph{(finite regions never leave the sector)} For every finite
  $R=[a,b]$, $\omega_\alpha\circ\Ad(U_R(g)^\dagger)\in\Kk_{\alpha\alpha}$;
  explicitly it is the state decorated with $A_\alpha$ outside $R$, $A_\beta$
  inside $R$, $V_\alpha(g)^{-1}$ on $\partial_-R$ and $V_\alpha(g)$ on
  $\partial_+R$ --- a kink at one edge and an antikink at the other, of total
  topological charge zero.
\item[(b)] \emph{(the half-infinite limit exists, weak-$*$ only, with a rate)}
  For $O\in\alg_W$, $w:=\max W$, $y>w$, and for \textbf{every}
  $\tilde\lambda\in(\lambda_E,1)$ there is $C_{\tilde\lambda}<\infty$ with
  \begin{equation}
    \Bigl|\,\omega_\alpha\bigl(U_{[x,y]}(g)^\dagger O\,U_{[x,y]}(g)\bigr)
      -\varrho^{(g)}_x(O)\,\Bigr|
    \;\le\; C_{\tilde\lambda}\,\norm{O}\,\tilde\lambda^{\,y-w},
    \label{eq:A2-rate}
  \end{equation}
  where $\varrho^{(g)}_x$ is the two-sided decorated state with $A_\alpha$ on
  $(-\infty,x-1]$, $V_\alpha(g)^{-1}$ on the bond $\bond{x-1}{x}$, and $A_\beta$
  on $[x,\infty)$.  Hence
  $1_{[x,\infty)}g\act\omega_\alpha=\varrho^{(g)}_x$ in the weak-$*$ topology.
\item[(c)] \emph{(the limit is a kink)} $\varrho^{(g)}_x\in\Kk_{\alpha\beta}$.
\item[(d)] \emph{(the sector jump)} Every finite-$R$ state is normal with respect
  to $\omega_\alpha$, while the limit is disjoint from it.  A weak-$*$ limit of a
  path of vector states thus leaves the folium of the vacuum; the jump is
  invisible at every finite $\abs{R}$ and is created entirely by the surviving
  end.
\item[(e)] \emph{(vacuum pairs, under hypothesis (T))} $G_L\times G_R$ acts
  componentwise on
  $\Omega_{\mathrm{vac}}\times\Omega_{\mathrm{vac}}$.  Assume \textbf{(T)}: $G$
  acts transitively on $\Omega_{\mathrm{vac}}$ (this does not follow from
  covariance --- the family may split into several $G$-orbits --- and without it
  every statement here holds per $G$-orbit).  Then
  \[
    \Omega_{\mathrm{vac}}\times\Omega_{\mathrm{vac}}
    \;\cong\;(G/H_\alpha)\times(G/H_\alpha),
  \]
  transitively, with stabiliser of $(\alpha',\beta')$ equal to
  $H_{\alpha'}\times H_{\beta'}$; and the complete invariant of a vacuum pair
  modulo the global diagonal symmetry is the \emph{double coset}
  \begin{equation}
    \mathfrak{d}(\alpha_L,\alpha_R)\;=\;H_\alpha\,g_L^{-1}g_R\,H_\alpha
    \;\in\;H_\alpha\backslash G/H_\alpha ,
    \label{eq:double-coset}
  \end{equation}
  which for $G=SU(2)$, $H_\alpha=\Ug(1)$ is the relative angle
  $\cos\vartheta=\hat n_L\cdot\hat n_R$.
\item[(f)] \emph{(every such sector is reached)} Every kink sector
  $\Kk_{\alpha\beta}$ with $\beta\in G\cdot\alpha$ arises this way, by choosing
  $g$ with $g\cdot\alpha=\beta$.
\end{enumerate}
\end{theorem}

\begin{scope}
The convergence in \cref{eq:A2-rate} is weak-$*$: the constant depends on the
observable's window through $w$, and the convergence cannot be improved to norm
convergence in $\alg^*$, because the approximants are normal with respect to the
vacuum and the limit is not.  The operators $U_{[x,y]}(g)$ have no strong limit
on the vacuum Hilbert space.  The rate is $\tilde\lambda\in(\lambda_E,1)$ and
\emph{not} $\lambda_E$ itself, which would be false if the transfer operator has
a Jordan block at the subleading modulus.  Hypothesis (T) is carried explicitly
and is not implied by covariance of the vacuum family.  Finally, the slogan ``a
kink is the contact term of a broken truncated symmetry'' is a mnemonic, not a
result: no contact term is defined anywhere in this campaign, no distributional
Ward identity is written, and no matrix element is shown to be supported at
coincident insertions.  It may not appear as the content of a proved claim.
\end{scope}

\provenance{A2}%
{\texttt{theory/corner-a-kinks.md}, steps (1.8)--(1.9)}%
{\texttt{theory/checks/corner\_a\_check.py} C7, C11}%
{\texttt{theory/verdicts/corner-a-r3.md}}

\smallskip\noindent\textit{Idea of proof.}
Clause (a) is \cref{thm:wi} read as a state: the decoration equals $A_\alpha$ on
both far tails, so both asymptotic labels are $\alpha$, and conjugation by a
local unitary preserves normality.  Clause (b) is a transfer-matrix gap estimate.
The finite-string state and the limit state carry identical decorations to the
left of the observable and differ only in the environment to its right, which is
$E_\beta^{\,y-w}\bigl(V_\alpha(g)r_\alpha V_\alpha(g)^\dagger\bigr)$ for the
first and $r_\beta$ for the second.  Since
$\operatorname{tr}(V_\alpha(g)r_\alpha V_\alpha(g)^\dagger)=1$, the difference is
controlled by the convergence $\norm{E_\beta^m(Y)-\operatorname{tr}(Y)r_\beta}\le
C_{\tilde\lambda}\tilde\lambda^m\norm{Y}$, valid for every
$\tilde\lambda>\lambda_E$.  One further point makes this uniform along the orbit:
by \cref{eq:IT} and unitarity of $u(g)$ one has
$E_\beta=\Ad_{V_\alpha(g)}\circ E_\alpha\circ\Ad_{V_\alpha(g)^{-1}}$, a
similarity, so the two transfer operators have the same spectrum and the same
gap.  Clauses (c), (d) then follow by reading the asymptotic labels and invoking
disjointness of distinct sectors.  For (e), orbit--stabiliser in each factor gives
the product structure; and while $g_L^{-1}g_R$ is invariant under the diagonal
action and changes to $h_1^{-1}(g_L^{-1}g_R)h_2$ under the residual stabilisers,
which is precisely the double-coset statement, the quantity $g_Lg_R^{-1}$
transforms by conjugation and is not diagonal-invariant at all.\hfill$\square$

\begin{remark}[Which model is which]
For the easy-axis chain the broken group is the discrete spin flip,
$\Omega_{\mathrm{vac}}=\{\uparrow,\downarrow\}$, and \cref{thm:a2} says that the
half-infinite flip applied to the polarised state produces a sharp domain wall
disjoint from the ferromagnetic vacuum: a genuinely gapped topological sector,
and the setting in which the superselection language is sharp
(\cref{sec:kink-model}).  For the isotropic ferromagnet $G=SU(2)$ is broken to
$\Ug(1)$ with $\Omega_{\mathrm{vac}}\cong S^2$ a continuum, $\chi=1$ and
$\lambda_E=0$, so the estimate \cref{eq:A2-rate} is exact at every $y>w$; but the
kinks interpolate between infinitesimally different vacua and cost arbitrarily
little energy.  The phrase ``the broken symmetry'' is ambiguous in this campaign
and is avoided: every statement names $G$, $H_\alpha$, or the broken complement
$\mathfrak{m}_\alpha$.
\end{remark}

\subsection{Goldstone tensors, the lattice Noether identity, and the kinematic factor}

The infinitesimal shadow of \cref{thm:wi} is a summation by parts.  The next
lemma is stated with its two boundary terms because they are $\Theta(1)$ and
dropping them is the error the round-one draft made.

\begin{lemma}[Summation by parts, with the exact gauge remainder]
\label{thm:lemma-sbp}
Let $\Lambda=[a,b]$ be finite, $X\in M_\chi(\C)$, and $f:\Z\to\C$ any profile.
Write $\ket{\psi;X@m}$ for the window vector with $X$ inserted on the bond
$\bond{m}{m+1}$, and $(\Delta f)(m):=f(m+1)-f(m)$.  Then, exactly,
\begin{equation}
  \sum_{n\in\Lambda} f(n)\,\bigl|\psi\bigl((A_\alpha X-XA_\alpha)@n\bigr)\bigr\rangle
  = -\sum_{m=a}^{b-1}(\Delta f)(m)\,\ket{\psi;X@m}
    + f(b)\,\ket{\psi;X@b} - f(a)\,\ket{\psi;X@(a-1)},
  \label{eq:SBP}
\end{equation}
the last two terms being the gauge remainder $\mathfrak{B}_\Lambda[f,X]$.  In
particular, for a plane wave $f(n)=e^{ikn}$ and the null map
$\mathcal{N}_k(X)^s:=e^{ik}A_\alpha^sX-XA_\alpha^s$,
\begin{equation}
  \bigl|\Phi^\Lambda_k(\mathcal{N}_k(X))\bigr\rangle
  = e^{ik(b+1)}\ket{\psi;X@b} - e^{ika}\ket{\psi;X@(a-1)}
  \;\neq\;0 .
  \label{eq:null-remainder}
\end{equation}
Moreover $\norm{\mathfrak{B}_\Lambda[f,X]}\le 2C_\partial\norm{X}\max(\abs{f(a)},\abs{f(b)})$
uniformly in $\Lambda$; if $f\in c_0(\Z)$ the remainder vanishes in norm as
$\Lambda\nearrow\Z$ (a statement about the remainder only), and if moreover
$f\in\ell^1\cap BV$ both sides are absolutely convergent and the identity is
remainderless in the limit; for a plane wave the remainder is $O(1)$, hence
$O(\abs{\Lambda}^{-1/2})$ after $\delta$-normalisation.  Finally
$\operatorname{rank}\mathcal{N}_k=\chi^2$ for $k\neq0$ and
$\operatorname{rank}\mathcal{N}_0=\chi^2-1$: the rank \emph{drops} by one as
$k\to0$.
\end{lemma}

\provenance{--- (supporting lemma of G0; no separate claim row)}%
{\texttt{theory/corner-a-goldstone.md}, step (1.5)}%
{\texttt{theory/checks/corner\_a\_check.py} C3, C3b, C4, C5, C10}%
{\texttt{theory/verdicts/corner-a-r3.md}}

\smallskip\noindent\textit{Idea of proof.}
Placing $A_\alpha^sX$ at site $n$ is the same as placing $X$ on the bond
immediately to its right, and placing $XA_\alpha^s$ at site $n$ is placing $X$ on
the bond to its left; subtracting the two finite sums after a shift of index
leaves the difference coefficient $-(\Delta f)(m)$ on the interior bonds and one
uncancelled term at each end.  The uniform bound is one insertion's worth of
norm, uniformly in the window because the transfer powers are power-bounded.
Numerically the identity holds to $3\cdot10^{-17}$, the plane-wave remainder has
norm $1.771$ for a $\chi=2$ Pauli tensor at $k=0.37$, and the ratio of remainder
to bulk for the centred profile $(1+\abs{n-c})^{-3}$ falls as
$3.4\cdot10^{-1},\,5.7\cdot10^{-2},\,8.5\cdot10^{-3},\,1.2\cdot10^{-3}$ for
window half-widths $4,8,16,32$.  No lower bound on the growth of the bulk term is
claimed or needed: for $\chi=1$ the plane-wave bulk sum is a bounded geometric
sum.  The rank statement is injectivity of $\mathcal{N}_k$ for $k\neq0$ (iterating
the relation forces $e^{ik}=1$) together with $\ker\mathcal{N}_0=\C\mathbf{1}$;
the factor $1-e^{ik}$ of \cref{thm:g0}(c) is algebraic and is \emph{not} caused by
this rank discontinuity.\hfill$\square$

\begin{theorem}[Goldstone tensors, the exact finite-window identity, and lattice Noether (G0)]
\label{thm:g0}
\statusProved\;
Let $\xi\in\mathfrak{g}$ and let $B_G(\xi)$ be the Goldstone tensor of
\cref{sec:symmetry-noether}, with $X_\alpha(\xi)$ the associated virtual
generator.
\begin{enumerate}
\item[(a)] \emph{(exact form on unbroken directions)} For $\xi\in\mathfrak{h}_\alpha$,
  $B_G(\xi)=A_\alpha X_\alpha(\xi)-X_\alpha(\xi)A_\alpha+i\theta'_\alpha(\xi)A_\alpha$,
  and after normal ordering the last term is absent, so
  $B_G(\xi)=\mathcal{N}_0(X_\alpha(\xi))$.
\item[(b)] \emph{(the dichotomy)}
  $B_G(\xi)\in\ran\mathcal{N}_0+\C A_\alpha$ \textbf{iff}
  $\xi\in\mathfrak{h}_\alpha$.  For a broken direction it is \emph{not} the case
  that the intertwining relation fails --- \cref{eq:IT} holds for every $g$, with
  target tensor $A_{g\cdot\alpha}$.  What fails is the \emph{same-vacuum return}
  $g\cdot\alpha=\alpha$, so $B_G(\xi)$ is a tangent to the vacuum manifold rather
  than to the gauge orbit.
\item[(c)] \emph{(the $k$-dependence, with the limit named)} For
  $\xi\in\mathfrak{h}_\alpha$, normally ordered, and any $k$,
  $B_G(\xi)=\mathcal{N}_k(X)+(1-e^{ik})A_\alpha X$ with $X:=X_\alpha(\xi)$; hence
  on a finite window $\Lambda=[a,b]$ the \textbf{exact} identity is
  \begin{equation}
    \bigl|\Phi^\Lambda_k(B_G(\xi))\bigr\rangle
    = (1-e^{ik})\sum_{m=a}^{b-1}e^{ikm}\ket{\psi;X@m}
      + e^{ikb}\ket{\psi;X@b} - e^{ika}\ket{\psi;X@(a-1)} .
    \label{eq:G0c}
  \end{equation}
  The clean form $(1-e^{ik})\bigl|\Phi_k(A_\alpha X_\alpha(\xi))\bigr\rangle$
  holds \textbf{only} in the $\delta$-normalised sense; for decaying profiles the
  correct statement is the real-space identity \cref{eq:SBP} with
  $f\in\ell^1\cap BV$, not a fixed-$k$ equation.
\item[(d)] \emph{(lattice Noether: the bond potential)} Exactly, on the vacuum,
  with no limit and no boundary term, for $\xi\in\mathfrak{h}_\alpha$ normally
  ordered,
  \begin{equation}
    q_x(\xi)\act\omega_\alpha
    \;=\;\bigl(\bondcur_{x|x+1}(\xi)-\bondcur_{x-1|x}(\xi)\bigr)\act\omega_\alpha:
    \label{eq:G0d}
  \end{equation}
  the \emph{physical} charge density acting on the vacuum is the lattice
  divergence of the \emph{virtual}, bond-supported quantity $\bondcur$.
\item[(e)] \emph{(continuity equation, finite range)} For a translation-invariant
  $G$-invariant Hamiltonian of finite range, with the cut current
  $j_{m|m+1}(\xi):=-[H,\sum_{y\le m}q_y(\xi)]\in\alg_{\mathrm{loc}}$, one has
  $[H,q_x(\xi)]=j_{x-1|x}(\xi)-j_{x|x+1}(\xi)$ and, for every admissible profile
  $f$ of compact support and \textbf{any} $\xi\in\mathfrak{g}$,
  \begin{equation}
    \bigl[H,\;Q[f;\xi]\bigr] \;=\; \sum_x (\Delta f)(x)\, j_{x|x+1}(\xi),
    \qquad\text{whence}\qquad
    [H,Q_k(\xi)]=(e^{ik}-1)\,J_k(\xi)
    \label{eq:G0e}
  \end{equation}
  in the wave-packet sense.
\end{enumerate}
\end{theorem}

\begin{scope}
Clauses (a), (b), (d) are asserted for $\xi\in\mathfrak{h}_\alpha$ and normally
ordered; clause (e) holds for any $\xi\in\mathfrak{g}$ and any finite-range
Hamiltonian.  The prefactor $(e^{ik}-1)$ in \cref{eq:G0e} is \emph{kinematic}: it
comes from the profile alone, it is convention-dependent (a kernel $e^{-ikx}$
gives $e^{-ik}-1$, and reversing the current's orientation flips the sign), and
its independence of the Hamiltonian is the tautology that a discrete difference
was factored out \emph{after} an $H$-dependent current had been defined.  It is
not a soft factor; see \cref{ref:g0-soft}.  Clauses (a)--(d) are statements about
the matrix-product vacuum and the linear ansatz space; only (e) is an operator
identity in the local algebra, true of the model itself whatever its ground state
and needing no matrix-product structure --- which is why (e), and nothing else in
Corner A, is the candidate seed for the soft-theorem corner
(\cref{sec:soft-index}).
\end{scope}

\provenance{G0}%
{\texttt{theory/corner-a-goldstone.md}, steps (1.5)--(1.7)}%
{\texttt{theory/checks/corner\_a\_check.py} C3, C3b, C4, C5, C9, C10; oracle O1 rederived at step (1.7)(2.3)}%
{\texttt{theory/verdicts/corner-a-r3.md}}

\smallskip\noindent\textit{Idea of proof.}
Clause (a) is \cref{eq:IT} differentiated at the identity along an unbroken
direction, where the target tensor returns to $A_\alpha$; the product rule
produces the commutator plus the phase derivative, which normal ordering removes.
For the forward direction of (b), the state map is constant on gauge orbits and
on rays, so its differential annihilates $\ran\mathcal{N}_0+\C A_\alpha$; the
resulting curve of vacua has vanishing derivative everywhere by the group law,
hence is constant, and since distinct labels give distinct states the direction
must stabilise $\alpha$.  Clause (c) is the algebraic decomposition
$\mathcal{N}_k(X)=(A_\alpha X-XA_\alpha)+(e^{ik}-1)A_\alpha X$ fed into
\cref{eq:SBP} and \cref{eq:null-remainder}, with care over the two summation
ranges: the bulk sum for $A_\alpha X$ runs to $b$, the null-map remainder
contributes at $b$ as well, and combining the two right-edge coefficients through
$(1-e^{ik})e^{ikb}+e^{ik(b+1)}=e^{ikb}$ gives \cref{eq:G0c}.  An earlier draft
truncated the bulk sum at $b-1$ while keeping the uncombined right coefficient,
thereby losing the term $(1-e^{ik})e^{ikb}\ket{\psi;X@b}$; the omission was
measured as $0.45058621$, exactly that term's norm, while both corrected displays
agree with the left-hand side to $6\cdot10^{-17}$.  Clause (d) is a single-site
statement --- the charge density replaces the tensor at $x$ by $B_G(\xi)$, i.e.\
by $X_\alpha(\xi)$ on the right bond minus $X_\alpha(\xi)$ on the left bond ---
so no reindexing and no boundary term arise; summing over a region telescopes it
back to the two insertions of \cref{thm:wi}.  Clause (e) is the observation that
the cut current is local, because $G$-invariance kills every term of the
Hamiltonian lying entirely inside the half-line and locality kills every term
lying entirely outside, leaving only the finitely many straddling terms; one Abel
summation then gives \cref{eq:G0e}.\hfill$\square$

\begin{proposition}[The isotropic ferromagnet, and the one oracle Corner A rederives]
\label{thm:fm}
\statusProved\;
For the isotropic ferromagnet with $G=SU(2)$ in the spin-$\tfrac12$
representation and the polarised product vacuum ($\chi=1$, $\lambda_E=0$):
(i) $H_\alpha=\Ug(1)$, the rotations about $z$; (ii) $[\omega_\alpha]=0$ and the
unbroken $\Ug(1)$ produces no Goldstone tensor; (iii) the two broken real
generators produce \emph{one} complex Goldstone tensor,
$B_G(\xi_y)=i\,B_G(\xi_x)$ --- the type-B count, derived rather than assumed;
(iv) $\ket{\Phi_k(B_G(\xi_x))}\propto\ket{k}$ is the exact one-magnon state and
\cref{eq:G0e} reproduces the dispersion $\omega(k)=J(1-\cos k)$, which is the
content of oracle fact O1.
\end{proposition}

\provenance{--- (supporting proposition of G0; no separate claim row)}%
{\texttt{theory/corner-a-goldstone.md}, step (1.7)}%
{\texttt{theory/checks/corner\_a\_check.py} C0; oracle O1}%
{\texttt{theory/verdicts/corner-a-r3.md}}

\smallskip\noindent\textit{Idea of proof.}
The $z$-rotation acts on the polarised site by a phase, so \cref{eq:IT} holds with
$V_\alpha=1$ and $\theta'_\alpha=\tfrac12\neq0$ --- the extensive phase of
\cref{thm:wi}, which is exactly the vacuum magnetisation density that normal
ordering subtracts.  Rotations off the axis send the tensor to an inequivalent
one, so the stabiliser is the circle; its second cohomology vanishes and
$V_\alpha\equiv1$, giving (ii).  For (iii), the raising direction annihilates the
vacuum site-wise and the lowering direction creates the magnon, so the real
two-dimensional broken space maps onto a one-dimensional complex space of
Goldstone tensors.  For (iv), the ferromagnetic current itself is a difference of
broken charge densities, supplying a \emph{second} factor of $(1-e^{ik})$; the two
factors together produce the quadratic dispersion.  This is the only oracle fact
that Corner A rederives, and only in the one-magnon sector.\hfill$\square$

\subsection{Two claims that were withdrawn}

Both of the following were asserted in earlier rounds of this corner, are false
as stated, and are recorded here because each killed a mechanism the campaign had
been relying on.

\begin{refutedclaim}[The naive asymptotic symmetry group is not the classifier (A2-orbit-r1)]
\label{ref:a2-orbit}
\normalfont
\statusRefuted\;
\emph{Withdrawn statement.}  The coset space $\A=(G_L\times G_R)/\Gdiag$ is the
lattice asymptotic symmetry group; its orbit is the set of vacuum pairs, and it
is labelled by the class $[g_Lg_R^{-1}]$.

\smallskip\noindent
\emph{Why it fails.}  Three independent reasons.  (1) $\A$ is a group only for
abelian $G$, since $\Gdiag$ is normal in $G\times G$ exactly then.  (2) The
componentwise stabiliser of a vacuum pair is $H_{\alpha}\times H_{\beta}$, not
$\Gdiag$, so the pair space is $(G/H_\alpha)\times(G/H_\alpha)$ and has the wrong
dimension to be $\A$: for $G=SU(2)$, $H_\alpha=\Ug(1)$ that is
$S^2\times S^2$, of dimension $4$, against $\dim SU(2)=3$.  (3) The proposed
label is not diagonal-invariant: under $g_i\mapsto hg_i$ the quantity
$g_Lg_R^{-1}$ transforms by conjugation, $g_Lg_R^{-1}\mapsto hg_Lg_R^{-1}h^{-1}$.
The numerical gate makes the contrast explicit: the relative angle is invariant to
$1.4\cdot10^{-16}$, while
$\norm{g_Lg_R^{-1}-h\,g_Lg_R^{-1}h^{-1}}=1.364$.  In the unbroken case the same
error appears in a different guise: the orbit is $\Aeff=G/N_\alpha$, and a
$\chi=1$ symmetric product vacuum has $N_\alpha=G$ and a \emph{one-point} orbit
where the withdrawn claim predicted an orbit of size $\abs{G}$.

\smallskip\noindent
\emph{What survives.}  \cref{thm:a1}(e) in the unbroken case: the orbit is
$\Aeff=G/N_\alpha$, a genuine group, with stabiliser
$S_\alpha\supseteq\Gdiag$.  And \cref{thm:a2}(e) in the broken case: the pair
space is $(G/H_\alpha)\times(G/H_\alpha)$ and the complete diagonal invariant is
the double coset \cref{eq:double-coset}.  Consequently the phrase ``asymptotic
symmetry group'' may not be used in this campaign without saying which of
$\Aeff$, $(G/H_\alpha)\times(G/H_\alpha)$, or $H_\alpha\backslash G/H_\alpha$ is
meant.
\end{refutedclaim}

\provenance{A2-orbit-r1}%
{disproved in \texttt{theory/corner-a-kinks.md}, step (1.9)(2.5)(3.3)}%
{\texttt{theory/checks/corner\_a\_check.py} C7}%
{\texttt{theory/verdicts/corner-a-r3.md}}

\begin{refutedclaim}[The naive universal soft factor (G0-soft-r1)]
\label{ref:g0-soft}
\normalfont
\statusRefuted\;
\emph{Withdrawn statement.}  The factor $(e^{ik}-1)$ appearing in
\cref{eq:G0e} is a \emph{universal soft factor}: it yields an Adler zero and
forces the hard-independent linear coefficient of the oracle's two-body phase.

\smallskip\noindent
\emph{Why it fails.}  In a matrix element the continuity equation says only
\begin{equation}
  \bra{\mathrm{out}}[H,Q_k]\ket{\mathrm{in}}
  =(e^{ik}-1)\,\bra{\mathrm{out}}J_k\ket{\mathrm{in}} .
  \label{eq:no-adler}
\end{equation}
\emph{All} dependence on the Hamiltonian, on the direction $\xi$, on the vacuum
and on the hard legs sits in the current matrix element on the right.  If
$\bra{\mathrm{out}}J_k\ket{\mathrm{in}}=C_{\mathrm{hard}}+O(k)$, then hard data
enters the amplitude at order $k$, not $k^2$, and there is no universal linear
coefficient; if that matrix element has a $1/k$ infrared singularity, there is no
zero at all.  The factor is moreover convention-dependent, and its independence of
the Hamiltonian is a tautology, as recorded in the scope note to \cref{thm:g0}.

\smallskip\noindent
\emph{What survives, and what it costs.}  \cref{thm:g0}(e) itself, as a lattice
continuity equation, exactly and for any finite-range Hamiltonian; and
\cref{thm:fm}, in which the ferromagnetic current happens to supply a second
lattice difference and the oracle's dispersion is genuinely rederived.  What is
lost is the brief's original mechanism for the soft theorem.  Recovering it is
the burden of the soft corner (\cref{sec:soft-index}, \cref{sec:universality}),
which must supply four things Corner A does not: regularity of
$\bra{\mathrm{out}}J_k\ket{\mathrm{in}}$ at $k=0$ on the relevant scattering
states; a reduction turning the operator identity \cref{eq:no-adler} into a
statement about amplitudes; a Ward identity relating the zero-momentum current to
the external legs' charges; and control of the remainder in a stated norm.  Until
those exist, agreement with the oracle's linear coefficient is a target, not a
derivation.
\end{refutedclaim}

\provenance{G0-soft-r1}%
{withdrawn in \texttt{theory/corner-a-goldstone.md}, step (1.6)(2.7)}%
{---}%
{\texttt{theory/verdicts/corner-a-r3.md}}

\subsection{Scope of the corner as a whole}

\begin{scope}[Four amendments, and two sketches that may not be used silently]
\label{scope:corner-a-sketches}
Four corrections must be absorbed by everything downstream of this section.
\begin{enumerate}
\item[(i)] \emph{The naive asymptotic symmetry group is not the classifying
  object.}  Unbroken: the orbit is $\Aeff=G/N_\alpha$, a group.  Broken, under
  hypothesis (T) --- otherwise per $G$-orbit --- vacuum pairs form
  $(G/H_\alpha)\times(G/H_\alpha)$ and the diagonal invariant is the double coset
  in $H_\alpha\backslash G/H_\alpha$.  The coset space $\A$ is a group only for
  abelian $G$ and is never the orbit unless $N_\alpha=\{e\}$
  (\cref{ref:a2-orbit}).
\item[(ii)] \emph{The edge from this corner to the soft corner is not supplied
  here.}  \cref{thm:g0}(e) is a lattice continuity equation; it gives no Adler
  zero, no universality, and does not rederive the two-body oracle facts
  (\cref{ref:g0-soft}).
\item[(iii)] \emph{``The zero-momentum Goldstone mode is pure gauge'' holds only
  along unbroken directions.}  Along a broken direction the intertwining relation
  \cref{eq:IT} still holds, with target tensor $A_{g\cdot\alpha}$; what fails is
  the same-vacuum return (\cref{thm:g0}(b)).
\item[(iv)] \emph{The ansatz gauge identity is not a finite-window identity.}  It
  carries two boundary terms of size $\Theta(1)$ (\cref{eq:null-remainder}) and
  holds only in the named limits: norm convergence for summable profiles of
  bounded variation, and $\delta$-normalisation for plane waves.
\end{enumerate}
Two sketches in this corner are load-bearing and may not be used silently.  The
first is the split property: it is not proved that a bond-decorated state
$\omega_\alpha^{M@b}$ is a normal state of the vacuum's von Neumann algebra,
equivalently that the charge algebra $\chargealg$ acts on the GNS Hilbert space
rather than on finite-window spaces.  The expected route is an infinite-volume
Schmidt decomposition across the bond, with the Schmidt index identified with the
matrix-product virtual index; at finite windows this is elementary, in the
thermodynamic limit it is open.  Its absence is exactly why the background
desideratum of a charge algebra acting on the physical state space is not met by
\cref{thm:a1}; every downstream use of ``the asymptotic charge algebra acting on
the vacuum Hilbert space'' must cite this gap.  It is \emph{not} needed for
\cref{thm:wi}, for clauses (a)--(c), (e)--(g) of \cref{thm:a1}, for
\cref{thm:a2}, or for \cref{thm:g0}.  The second is uniformity over a continuous
vacuum manifold: for the isotropic ferromagnet, \cref{thm:a2} produces an
uncountable family of mutually disjoint sectors with no uniform separation --- the
separating observable degrades as $g\to e$, so the sectors are algebraically
superselected but not energetically separated, the kink creation energy tending to
zero.  Missing are a version of \cref{thm:a2} uniform in $g$, with the separation
quantified by the distance between the two vacua, and a selection criterion that
cuts the continuum down to a physically meaningful family --- exactly the question
of which function spaces are admissible at infinity.  Consequently the
ferromagnet's kinks may not be used as superselection sectors in the soft-theorem
conjecture, and the memory construction of \cref{sec:memory-quantization} cannot
be run on that model through this route.  Nothing in \cref{thm:a2}(a)--(f) is
affected: those are proved for each fixed $g\notin H_\alpha$.
\end{scope}

\provenance{WI, A1, A2, G0 (scope common to all four)}%
{\texttt{theory/corner-a-pitfalls.md} (campaign pitfall list); sketches at \texttt{theory/corner-a.md} step (1.4)(2.9) and \texttt{theory/corner-a-kinks.md} step (1.10)(2.3)}%
{\texttt{theory/checks/corner\_a\_check.py} C0--C11 (all pass)}%
{\texttt{theory/verdicts/corner-a-r3.md} (loop converged; four minor wording notes only)}
